\subsection{Reconstruction}
A segmentation output has, for each camera and for each frame of a sequence, a set of coordinates $(x, y)$ that locate the position in the image of those markers that were detected.
The reconstruction process is to obtain the position of the markers in space, from the position of the markers in at least two retinas.

The reconstruction process presented was inspired by the work of Herda and consists of three basic steps:
\vspace{-0.2cm}
\begin{enumerate}
\item Find the correspondence between points in different retinas.
\item Select the best match.
\item Reconstruct and check the rest of the retinas.
\end{enumerate}
\vspace{-0.8cm}
\subsubsection{Algorithm}
The implemented algorithm receives as input 2D points of the detected markers and returns as output reconstructed 3D points.
Figure \ref{fig: diagrama algoritmo} is a diagram of the algorithm presented.
%\vspace{-0.2cm}
\begin{figure}
    \begin{center}
        \includegraphics[scale=0.45]{./imagenes/Reconstruccion/bloques_reconstruccion}
        \caption{Diagrama de bloques del algoritmo de reconstrucción.}
        \label{fig: diagrama algoritmo}
    \end{center}
\end{figure}
%\vspace{-0.6cm}

\textit{Associating 2D points.}\label{seccion_asociar2D_uno}
%
This block receives as input the coordinates of the detected points and projection matrices of the cameras, returning to each point a list sorted by relevance, existing partnerships with points on other cameras. \\ 
\textit{\hspace*{0.5cm}Best match.}\label{MejorAsociacion}
%
From the list of associations between pairs of cameras it is necessary to choose one that possesses more likely to form the pair corresponding to the projection of a 3D marker views on such images.
%
On each pair of cameras is taken that association which maintains the shorter distance and contains valid points, discarding the rest. 
\vspace{-0.65cm}
\begin{figure}
    \begin{center}
       \includegraphics[scale=1.0]{./imagenes/Reconstruccion/Algoritmo_reconstruccion}       
    \end{center}
\end{figure}
\vspace{-1.0cm}

For the choice of pairs of cameras two cases were considered.
The first evaluates each camera with all the remaining and the second considers the arrangement of the cameras in space and match the adjacent cameras consecutively.

\textit{3D reconstruction and Validation.}\label{seccion_reconstruccion3D_validacion}
$x_i x_j$ couple of points of the cameras $i$ and $j$ respectively,  reconstruct a valid point 3D $X_{ij}$ if there is at least one $x_k$ in camera $k\not= i, \,j$ such that $X_{ik} \in $ sphere$(X_{ij}, \delta)$, to certain threshold $\delta$.

\textit{Update matches}\label{actualizar_asociaciones}
The couple who reconstructs $X_{ij}$ as well  as the $x_k$ points who validate this reconstruction are removed. 
Iteration is continued repeating the process with the next pairs cameras best associated.
Finally the iterative process stops when the number of reconstructed markers equals the number of markers that person has, equal to the maximum number of  reconstructions indicated, or no valid 2D points such that a partnership between different points of view can be established.